\subsection{Memoria profunda y compatibilidad de los cilindros}
\label{memprof:seccion}

La prolongación compara estados que conservan simultáneamente residuos,
cocientes y fronteras. En la coordenatización lineal de seis coordenadas del registro,
tomamos el representante entero de \(L_0\) cuyas entradas se han mostrado:
\[
 A_H=\begin{pmatrix}
 2&2&2&1&2&1\\2&1&2&2&1&1\\1&0&0&1&0&2\\
 0&0&1&1&1&0\\2&2&2&0&2&1\\2&1&2&1&0&0
 \end{pmatrix},\qquad\det A_H=1.
\]
La elección de este representante forma parte de la coordenatización: la reducción
módulo tres conserva \(L_0\), mientras que el representante entero fija
la operación sobre los cocientes profundos. Con vectores fila se obtiene
\begin{equation}
 \begin{aligned}
 \mathsf{Div}_H:\mathbb Z^6&\longrightarrow
       \{0,1,2\}^6\times\mathbb Z^6,\\
 x_H&\longmapsto(u_H,x_H^+),\\
 y_H&=x_HA_H,\qquad u_H=[y_H]_3,\qquad
 x_H^+=(y_H-u_H)/3 .
 \end{aligned}
 \label{memprof:division}
\end{equation}
El par \((u_H,x_H^+)\) conserva \(y_H=u_H+3x_H^+\). Puesto que
\(A_H\) es unimodular, su inversa es entera y
\begin{equation}
 x_H=(u_H+3x_H^+)A_H^{-1}.
 \label{memprof:inversa}
\end{equation}
Esta identidad recupera la coordenada profunda desde su residuo y su
cociente. Los otros campos del estado conservan sus registros propios.
La orientación ternaria se lee después mediante
\(0\mapsto0,\ 1\mapsto+1,\ 2\mapsto-1\), manteniendo el residuo
\(u_H\in\{0,1,2\}^6\) y el cociente de \eqref{memprof:division}.

La coordenatización cilíndrica relaciona un prefijo ternario de longitud \(L\), de
índice entero \(N\), con \(k_{\mathrm{dec}}\) tríadas decimales, de
índice \(D\). Los cilindros semiabiertos correspondientes son
\[
 I_3(N,L)=\left[\frac N{3^L},\frac{N+1}{3^L}\right),\qquad
 I_{1000}(D,k_{\mathrm{dec}})
 =\left[\frac D{1000^{k_{\mathrm{dec}}}},
         \frac{D+1}{1000^{k_{\mathrm{dec}}}}\right).
\]
Definimos los enteros de frontera
\[
 a_{\mathrm{cyl}}=N1000^{k_{\mathrm{dec}}}-D3^L,\qquad
 b_{\mathrm{cyl}}=(N+1)1000^{k_{\mathrm{dec}}}-D3^L.
\]
Para un bloque de seis residuos fijamos su codificación entera
\[
 \nu(u_H)=\sum_{j=1}^6(u_H)_j3^{6-j}
 \in\{0,\ldots,728\}.
\]

\begin{lemma}[Actualización entera de los cilindros compatibles]
\label{memprof:cilindros}
Los cilindros anteriores se intersectan si y sólo si
\(a_{\mathrm{cyl}}<3^L\) y \(b_{\mathrm{cyl}}>0\).
Al añadir un bloque ternario de índice
\(\nu\in\{0,\ldots,728\}\), se tiene \(N'=729N+\nu\) y
\(L'=L+6\). Para \(\Delta k_{\mathrm{dec}}\in\{0,1\}\),
las fronteras se actualizan por
\[
 \begin{array}{ll}
 \Delta k_{\mathrm{dec}}=0:
   &a_{\mathrm{cyl}}'=729a_{\mathrm{cyl}}+\nu1000^{k_{\mathrm{dec}}},\\[2pt]
 \Delta k_{\mathrm{dec}}=1:
   &D'=1000D+\delta,\\
   &a_{\mathrm{cyl}}'=729000a_{\mathrm{cyl}}
      +\nu1000^{k_{\mathrm{dec}}+1}-\delta\,729\,3^L ,
 \end{array}
\]
donde \(\delta\in\{0,\ldots,999\}\) es la tríada de la prolongación
compatible. En ambos casos,
\[
 b_{\mathrm{cyl}}'=a_{\mathrm{cyl}}'
          +1000^{k_{\mathrm{dec}}+\Delta k_{\mathrm{dec}}}.
\]
\end{lemma}
\begin{proof}
Dos intervalos semiabiertos se intersectan precisamente cuando cada extremo
izquierdo es menor que el extremo derecho del otro. Multiplicar esas dos
desigualdades por \(3^L1000^{k_{\mathrm{dec}}}\) da
\(a_{\mathrm{cyl}}<3^L\) y \(b_{\mathrm{cyl}}>0\).
La sustitución de \(N'\), \(L'\) y \(D'\) en
\[
 a_{\mathrm{cyl}}'
 =N'1000^{k_{\mathrm{dec}}+\Delta k_{\mathrm{dec}}}-D'3^{L'}
\]
produce las dos recurrencias. Restar las definiciones de
\(b_{\mathrm{cyl}}'\) y \(a_{\mathrm{cyl}}'\) da la última identidad.
\end{proof}

Estas operaciones actúan sobre prefijos enteros y sus restricciones
compatibles. La lectura arquimediana aparece al evaluar el sistema de
cilindros resultante. El dominio de prolongación incorpora además la
orientación, el acarreo, la firma conjunta y los datos de frontera del estado
\(\widehat X\) de \eqref{eq:actualizacion}; la comparación de intervalos
es una componente de esa compatibilidad conjunta.
